% Options for packages loaded elsewhere
\PassOptionsToPackage{unicode}{hyperref}
\PassOptionsToPackage{hyphens}{url}
\documentclass[
]{article}
\usepackage{xcolor}
\usepackage[margin=1in]{geometry}
\usepackage{amsmath,amssymb}
\setcounter{secnumdepth}{-\maxdimen} % remove section numbering
\usepackage{iftex}
\ifPDFTeX
  \usepackage[T1]{fontenc}
  \usepackage[utf8]{inputenc}
  \usepackage{textcomp} % provide euro and other symbols
\else % if luatex or xetex
  \usepackage{unicode-math} % this also loads fontspec
  \defaultfontfeatures{Scale=MatchLowercase}
  \defaultfontfeatures[\rmfamily]{Ligatures=TeX,Scale=1}
\fi
\usepackage{lmodern}
\ifPDFTeX\else
  % xetex/luatex font selection
\fi
% Use upquote if available, for straight quotes in verbatim environments
\IfFileExists{upquote.sty}{\usepackage{upquote}}{}
\IfFileExists{microtype.sty}{% use microtype if available
  \usepackage[]{microtype}
  \UseMicrotypeSet[protrusion]{basicmath} % disable protrusion for tt fonts
}{}
\makeatletter
\@ifundefined{KOMAClassName}{% if non-KOMA class
  \IfFileExists{parskip.sty}{%
    \usepackage{parskip}
  }{% else
    \setlength{\parindent}{0pt}
    \setlength{\parskip}{6pt plus 2pt minus 1pt}}
}{% if KOMA class
  \KOMAoptions{parskip=half}}
\makeatother
\usepackage{color}
\usepackage{fancyvrb}
\newcommand{\VerbBar}{|}
\newcommand{\VERB}{\Verb[commandchars=\\\{\}]}
\DefineVerbatimEnvironment{Highlighting}{Verbatim}{commandchars=\\\{\}}
% Add ',fontsize=\small' for more characters per line
\usepackage{framed}
\definecolor{shadecolor}{RGB}{248,248,248}
\newenvironment{Shaded}{\begin{snugshade}}{\end{snugshade}}
\newcommand{\AlertTok}[1]{\textcolor[rgb]{0.94,0.16,0.16}{#1}}
\newcommand{\AnnotationTok}[1]{\textcolor[rgb]{0.56,0.35,0.01}{\textbf{\textit{#1}}}}
\newcommand{\AttributeTok}[1]{\textcolor[rgb]{0.13,0.29,0.53}{#1}}
\newcommand{\BaseNTok}[1]{\textcolor[rgb]{0.00,0.00,0.81}{#1}}
\newcommand{\BuiltInTok}[1]{#1}
\newcommand{\CharTok}[1]{\textcolor[rgb]{0.31,0.60,0.02}{#1}}
\newcommand{\CommentTok}[1]{\textcolor[rgb]{0.56,0.35,0.01}{\textit{#1}}}
\newcommand{\CommentVarTok}[1]{\textcolor[rgb]{0.56,0.35,0.01}{\textbf{\textit{#1}}}}
\newcommand{\ConstantTok}[1]{\textcolor[rgb]{0.56,0.35,0.01}{#1}}
\newcommand{\ControlFlowTok}[1]{\textcolor[rgb]{0.13,0.29,0.53}{\textbf{#1}}}
\newcommand{\DataTypeTok}[1]{\textcolor[rgb]{0.13,0.29,0.53}{#1}}
\newcommand{\DecValTok}[1]{\textcolor[rgb]{0.00,0.00,0.81}{#1}}
\newcommand{\DocumentationTok}[1]{\textcolor[rgb]{0.56,0.35,0.01}{\textbf{\textit{#1}}}}
\newcommand{\ErrorTok}[1]{\textcolor[rgb]{0.64,0.00,0.00}{\textbf{#1}}}
\newcommand{\ExtensionTok}[1]{#1}
\newcommand{\FloatTok}[1]{\textcolor[rgb]{0.00,0.00,0.81}{#1}}
\newcommand{\FunctionTok}[1]{\textcolor[rgb]{0.13,0.29,0.53}{\textbf{#1}}}
\newcommand{\ImportTok}[1]{#1}
\newcommand{\InformationTok}[1]{\textcolor[rgb]{0.56,0.35,0.01}{\textbf{\textit{#1}}}}
\newcommand{\KeywordTok}[1]{\textcolor[rgb]{0.13,0.29,0.53}{\textbf{#1}}}
\newcommand{\NormalTok}[1]{#1}
\newcommand{\OperatorTok}[1]{\textcolor[rgb]{0.81,0.36,0.00}{\textbf{#1}}}
\newcommand{\OtherTok}[1]{\textcolor[rgb]{0.56,0.35,0.01}{#1}}
\newcommand{\PreprocessorTok}[1]{\textcolor[rgb]{0.56,0.35,0.01}{\textit{#1}}}
\newcommand{\RegionMarkerTok}[1]{#1}
\newcommand{\SpecialCharTok}[1]{\textcolor[rgb]{0.81,0.36,0.00}{\textbf{#1}}}
\newcommand{\SpecialStringTok}[1]{\textcolor[rgb]{0.31,0.60,0.02}{#1}}
\newcommand{\StringTok}[1]{\textcolor[rgb]{0.31,0.60,0.02}{#1}}
\newcommand{\VariableTok}[1]{\textcolor[rgb]{0.00,0.00,0.00}{#1}}
\newcommand{\VerbatimStringTok}[1]{\textcolor[rgb]{0.31,0.60,0.02}{#1}}
\newcommand{\WarningTok}[1]{\textcolor[rgb]{0.56,0.35,0.01}{\textbf{\textit{#1}}}}
\usepackage{graphicx}
\makeatletter
\newsavebox\pandoc@box
\newcommand*\pandocbounded[1]{% scales image to fit in text height/width
  \sbox\pandoc@box{#1}%
  \Gscale@div\@tempa{\textheight}{\dimexpr\ht\pandoc@box+\dp\pandoc@box\relax}%
  \Gscale@div\@tempb{\linewidth}{\wd\pandoc@box}%
  \ifdim\@tempb\p@<\@tempa\p@\let\@tempa\@tempb\fi% select the smaller of both
  \ifdim\@tempa\p@<\p@\scalebox{\@tempa}{\usebox\pandoc@box}%
  \else\usebox{\pandoc@box}%
  \fi%
}
% Set default figure placement to htbp
\def\fps@figure{htbp}
\makeatother
\setlength{\emergencystretch}{3em} % prevent overfull lines
\providecommand{\tightlist}{%
  \setlength{\itemsep}{0pt}\setlength{\parskip}{0pt}}
\usepackage{bookmark}
\IfFileExists{xurl.sty}{\usepackage{xurl}}{} % add URL line breaks if available
\urlstyle{same}
\hypersetup{
  pdftitle={Extrair valores do raster nos pontos de rendimento hídrico},
  pdfauthor={Eileen Acosta},
  hidelinks,
  pdfcreator={LaTeX via pandoc}}

\title{Extrair valores do raster nos pontos de rendimento hídrico}
\author{Eileen Acosta}
\date{2026-05-19}

\begin{document}
\maketitle

Ler o raster Extrair QF\_v2 do raster para os pontos Ler a tabela
rendimento\_hidrologico\_mm\_ano.csv Juntar essa tabela com os pontos
usando Codigo Calcular a média das colunas mm\_\ldots{} Fazer o gráfico
1:1

\begin{Shaded}
\begin{Highlighting}[]
\CommentTok{\# ============================================================}
\CommentTok{\# PARTE 2 — LER RASTER E PONTOS}
\CommentTok{\# ============================================================}

\NormalTok{raster\_qf }\OtherTok{\textless{}{-}} \FunctionTok{rast}\NormalTok{(arquivo\_raster)}

\NormalTok{pontos }\OtherTok{\textless{}{-}} \FunctionTok{st\_read}\NormalTok{(arquivo\_pontos)}
\end{Highlighting}
\end{Shaded}

\begin{verbatim}
## Reading layer `estacoes_check' from data source 
##   `C:\trab\estacoes\estacoes_check.shp' using driver `ESRI Shapefile'
## Simple feature collection with 107 features and 4 fields
## Geometry type: POINT
## Dimension:     XY
## Bounding box:  xmin: -72.74122 ymin: -18.7136 xmax: -42.6956 ymax: 3.4381
## Geodetic CRS:  WGS 84
\end{verbatim}

\begin{Shaded}
\begin{Highlighting}[]
\FunctionTok{cat}\NormalTok{(}\StringTok{"}\SpecialCharTok{\textbackslash{}n}\StringTok{CRS do raster:}\SpecialCharTok{\textbackslash{}n}\StringTok{"}\NormalTok{)}
\end{Highlighting}
\end{Shaded}

\begin{verbatim}
## 
## CRS do raster:
\end{verbatim}

\begin{Shaded}
\begin{Highlighting}[]
\FunctionTok{print}\NormalTok{(}\FunctionTok{crs}\NormalTok{(raster\_qf))}
\end{Highlighting}
\end{Shaded}

\begin{verbatim}
## [1] "PROJCRS[\"SIRGAS 2000 / Brazil Albers\",\n    BASEGEOGCRS[\"SIRGAS 2000\",\n        DATUM[\"Sistema de Referencia Geocentrico para las AmericaS 2000\",\n            ELLIPSOID[\"GRS 1980\",6378137,298.257222101,\n                LENGTHUNIT[\"metre\",1]]],\n        PRIMEM[\"Greenwich\",0,\n            ANGLEUNIT[\"degree\",0.0174532925199433]],\n        ID[\"EPSG\",4674]],\n    CONVERSION[\"Brazil Albers\",\n        METHOD[\"Albers Equal Area\",\n            ID[\"EPSG\",9822]],\n        PARAMETER[\"Latitude of false origin\",-12,\n            ANGLEUNIT[\"degree\",0.0174532925199433],\n            ID[\"EPSG\",8821]],\n        PARAMETER[\"Longitude of false origin\",-54,\n            ANGLEUNIT[\"degree\",0.0174532925199433],\n            ID[\"EPSG\",8822]],\n        PARAMETER[\"Latitude of 1st standard parallel\",-2,\n            ANGLEUNIT[\"degree\",0.0174532925199433],\n            ID[\"EPSG\",8823]],\n        PARAMETER[\"Latitude of 2nd standard parallel\",-22,\n            ANGLEUNIT[\"degree\",0.0174532925199433],\n            ID[\"EPSG\",8824]],\n        PARAMETER[\"Easting at false origin\",5000000,\n            LENGTHUNIT[\"metre\",1],\n            ID[\"EPSG\",8826]],\n        PARAMETER[\"Northing at false origin\",10000000,\n            LENGTHUNIT[\"metre\",1],\n            ID[\"EPSG\",8827]]],\n    CS[Cartesian,2],\n        AXIS[\"(E)\",east,\n            ORDER[1],\n            LENGTHUNIT[\"metre\",1]],\n        AXIS[\"(N)\",north,\n            ORDER[2],\n            LENGTHUNIT[\"metre\",1]],\n    USAGE[\n        SCOPE[\"Statistical analysis.\"],\n        AREA[\"Brazil - onshore.\"],\n        BBOX[-33.78,-74.01,5.28,-34.74]],\n    ID[\"EPSG\",10857]]"
\end{verbatim}

\begin{Shaded}
\begin{Highlighting}[]
\FunctionTok{cat}\NormalTok{(}\StringTok{"}\SpecialCharTok{\textbackslash{}n}\StringTok{CRS dos pontos:}\SpecialCharTok{\textbackslash{}n}\StringTok{"}\NormalTok{)}
\end{Highlighting}
\end{Shaded}

\begin{verbatim}
## 
## CRS dos pontos:
\end{verbatim}

\begin{Shaded}
\begin{Highlighting}[]
\FunctionTok{print}\NormalTok{(}\FunctionTok{st\_crs}\NormalTok{(pontos))}
\end{Highlighting}
\end{Shaded}

\begin{verbatim}
## Coordinate Reference System:
##   User input: WGS 84 
##   wkt:
## GEOGCRS["WGS 84",
##     DATUM["World Geodetic System 1984",
##         ELLIPSOID["WGS 84",6378137,298.257223563,
##             LENGTHUNIT["metre",1]]],
##     PRIMEM["Greenwich",0,
##         ANGLEUNIT["degree",0.0174532925199433]],
##     CS[ellipsoidal,2],
##         AXIS["latitude",north,
##             ORDER[1],
##             ANGLEUNIT["degree",0.0174532925199433]],
##         AXIS["longitude",east,
##             ORDER[2],
##             ANGLEUNIT["degree",0.0174532925199433]],
##     ID["EPSG",4326]]
\end{verbatim}

\begin{Shaded}
\begin{Highlighting}[]
\FunctionTok{cat}\NormalTok{(}\StringTok{"}\SpecialCharTok{\textbackslash{}n}\StringTok{Número de pontos:"}\NormalTok{, }\FunctionTok{nrow}\NormalTok{(pontos), }\StringTok{"}\SpecialCharTok{\textbackslash{}n}\StringTok{"}\NormalTok{)}
\end{Highlighting}
\end{Shaded}

\begin{verbatim}
## 
## Número de pontos: 107
\end{verbatim}

\begin{Shaded}
\begin{Highlighting}[]
\FunctionTok{names}\NormalTok{(pontos)}
\end{Highlighting}
\end{Shaded}

\begin{verbatim}
## [1] "Codigo"   "Nome"     "Latitud"  "Longitd"  "geometry"
\end{verbatim}

\begin{Shaded}
\begin{Highlighting}[]
\CommentTok{\# ============================================================}
\CommentTok{\# PARTE 3 — AJUSTAR CRS DOS PONTOS PARA O CRS DO RASTER}
\CommentTok{\# ============================================================}

\CommentTok{\# O raster e os pontos precisam estar no mesmo sistema de coordenadas.}
\CommentTok{\# Aqui transformamos os pontos para o CRS do raster.}

\NormalTok{pontos\_raster\_crs }\OtherTok{\textless{}{-}} \FunctionTok{st\_transform}\NormalTok{(}
\NormalTok{  pontos,}
  \FunctionTok{crs}\NormalTok{(raster\_qf)}
\NormalTok{)}

\FunctionTok{cat}\NormalTok{(}\StringTok{"}\SpecialCharTok{\textbackslash{}n}\StringTok{CRS dos pontos após transformação:}\SpecialCharTok{\textbackslash{}n}\StringTok{"}\NormalTok{)}
\end{Highlighting}
\end{Shaded}

\begin{verbatim}
## 
## CRS dos pontos após transformação:
\end{verbatim}

\begin{Shaded}
\begin{Highlighting}[]
\FunctionTok{print}\NormalTok{(}\FunctionTok{st\_crs}\NormalTok{(pontos\_raster\_crs))}
\end{Highlighting}
\end{Shaded}

\begin{verbatim}
## Coordinate Reference System:
##   User input: PROJCRS["SIRGAS 2000 / Brazil Albers",
##     BASEGEOGCRS["SIRGAS 2000",
##         DATUM["Sistema de Referencia Geocentrico para las AmericaS 2000",
##             ELLIPSOID["GRS 1980",6378137,298.257222101,
##                 LENGTHUNIT["metre",1]]],
##         PRIMEM["Greenwich",0,
##             ANGLEUNIT["degree",0.0174532925199433]],
##         ID["EPSG",4674]],
##     CONVERSION["Brazil Albers",
##         METHOD["Albers Equal Area",
##             ID["EPSG",9822]],
##         PARAMETER["Latitude of false origin",-12,
##             ANGLEUNIT["degree",0.0174532925199433],
##             ID["EPSG",8821]],
##         PARAMETER["Longitude of false origin",-54,
##             ANGLEUNIT["degree",0.0174532925199433],
##             ID["EPSG",8822]],
##         PARAMETER["Latitude of 1st standard parallel",-2,
##             ANGLEUNIT["degree",0.0174532925199433],
##             ID["EPSG",8823]],
##         PARAMETER["Latitude of 2nd standard parallel",-22,
##             ANGLEUNIT["degree",0.0174532925199433],
##             ID["EPSG",8824]],
##         PARAMETER["Easting at false origin",5000000,
##             LENGTHUNIT["metre",1],
##             ID["EPSG",8826]],
##         PARAMETER["Northing at false origin",10000000,
##             LENGTHUNIT["metre",1],
##             ID["EPSG",8827]]],
##     CS[Cartesian,2],
##         AXIS["(E)",east,
##             ORDER[1],
##             LENGTHUNIT["metre",1]],
##         AXIS["(N)",north,
##             ORDER[2],
##             LENGTHUNIT["metre",1]],
##     USAGE[
##         SCOPE["Statistical analysis."],
##         AREA["Brazil - onshore."],
##         BBOX[-33.78,-74.01,5.28,-34.74]],
##     ID["EPSG",10857]] 
##   wkt:
## PROJCRS["SIRGAS 2000 / Brazil Albers",
##     BASEGEOGCRS["SIRGAS 2000",
##         DATUM["Sistema de Referencia Geocentrico para las AmericaS 2000",
##             ELLIPSOID["GRS 1980",6378137,298.257222101,
##                 LENGTHUNIT["metre",1]]],
##         PRIMEM["Greenwich",0,
##             ANGLEUNIT["degree",0.0174532925199433]],
##         ID["EPSG",4674]],
##     CONVERSION["Brazil Albers",
##         METHOD["Albers Equal Area",
##             ID["EPSG",9822]],
##         PARAMETER["Latitude of false origin",-12,
##             ANGLEUNIT["degree",0.0174532925199433],
##             ID["EPSG",8821]],
##         PARAMETER["Longitude of false origin",-54,
##             ANGLEUNIT["degree",0.0174532925199433],
##             ID["EPSG",8822]],
##         PARAMETER["Latitude of 1st standard parallel",-2,
##             ANGLEUNIT["degree",0.0174532925199433],
##             ID["EPSG",8823]],
##         PARAMETER["Latitude of 2nd standard parallel",-22,
##             ANGLEUNIT["degree",0.0174532925199433],
##             ID["EPSG",8824]],
##         PARAMETER["Easting at false origin",5000000,
##             LENGTHUNIT["metre",1],
##             ID["EPSG",8826]],
##         PARAMETER["Northing at false origin",10000000,
##             LENGTHUNIT["metre",1],
##             ID["EPSG",8827]]],
##     CS[Cartesian,2],
##         AXIS["(E)",east,
##             ORDER[1],
##             LENGTHUNIT["metre",1]],
##         AXIS["(N)",north,
##             ORDER[2],
##             LENGTHUNIT["metre",1]],
##     USAGE[
##         SCOPE["Statistical analysis."],
##         AREA["Brazil - onshore."],
##         BBOX[-33.78,-74.01,5.28,-34.74]],
##     ID["EPSG",10857]]
\end{verbatim}

\begin{Shaded}
\begin{Highlighting}[]
\CommentTok{\# ============================================================}
\CommentTok{\# PARTE 4 — EXTRAIR VALORES DO RASTER NOS PONTOS}
\CommentTok{\# ============================================================}

\CommentTok{\# Converter sf para SpatVector, formato usado pelo pacote terra}
\NormalTok{pontos\_vect }\OtherTok{\textless{}{-}} \FunctionTok{vect}\NormalTok{(pontos\_raster\_crs)}

\CommentTok{\# Extrair valores do raster}
\NormalTok{valores\_extraidos }\OtherTok{\textless{}{-}}\NormalTok{ terra}\SpecialCharTok{::}\FunctionTok{extract}\NormalTok{(}
\NormalTok{  raster\_qf,}
\NormalTok{  pontos\_vect}
\NormalTok{)}

\FunctionTok{head}\NormalTok{(valores\_extraidos)}
\end{Highlighting}
\end{Shaded}

\begin{verbatim}
##   ID    QF_v2
## 1  1 1461.153
## 2  2 2131.965
## 3  3 2834.586
## 4  4 2789.684
## 5  5 1095.260
## 6  6 2078.195
\end{verbatim}

\begin{Shaded}
\begin{Highlighting}[]
\FunctionTok{names}\NormalTok{(valores\_extraidos)}
\end{Highlighting}
\end{Shaded}

\begin{verbatim}
## [1] "ID"    "QF_v2"
\end{verbatim}

\begin{Shaded}
\begin{Highlighting}[]
\CommentTok{\# ============================================================}
\CommentTok{\# PARTE 5 — JUNTAR VALORES EXTRAÍDOS COM OS ATRIBUTOS DOS PONTOS}
\CommentTok{\# ============================================================}

\NormalTok{nome\_coluna\_raster }\OtherTok{\textless{}{-}} \FunctionTok{names}\NormalTok{(valores\_extraidos)[}\DecValTok{2}\NormalTok{]}

\NormalTok{valores\_extraidos }\OtherTok{\textless{}{-}}\NormalTok{ valores\_extraidos }\SpecialCharTok{\%\textgreater{}\%}
  \FunctionTok{rename}\NormalTok{(}
    \AttributeTok{QF\_v2 =} \FunctionTok{all\_of}\NormalTok{(nome\_coluna\_raster)}
\NormalTok{  )}

\NormalTok{pontos\_com\_id }\OtherTok{\textless{}{-}}\NormalTok{ pontos }\SpecialCharTok{\%\textgreater{}\%}
  \FunctionTok{mutate}\NormalTok{(}
    \AttributeTok{ID =} \FunctionTok{row\_number}\NormalTok{(),}
    \AttributeTok{Codigo =} \FunctionTok{as.character}\NormalTok{(Codigo)}
\NormalTok{  )}

\NormalTok{tabela\_final }\OtherTok{\textless{}{-}}\NormalTok{ pontos\_com\_id }\SpecialCharTok{\%\textgreater{}\%}
  \FunctionTok{st\_drop\_geometry}\NormalTok{() }\SpecialCharTok{\%\textgreater{}\%}
  \FunctionTok{left\_join}\NormalTok{(}
\NormalTok{    valores\_extraidos,}
    \AttributeTok{by =} \StringTok{"ID"}
\NormalTok{  ) }\SpecialCharTok{\%\textgreater{}\%}
  \FunctionTok{select}\NormalTok{(}
    \SpecialCharTok{{-}}\NormalTok{ID}
\NormalTok{  )}

\FunctionTok{head}\NormalTok{(tabela\_final)}
\end{Highlighting}
\end{Shaded}

\begin{verbatim}
##     Codigo                         Nome Latitud  Longitd    QF_v2
## 1 11200000                     TERESINA -4.3578 -69.7336 1461.153
## 2 11400000      SÃO PAULO DE OLIVENÇA -3.4569 -68.9119 2131.965
## 3 11500000 SANTO ANTÔNIO DO IÇÃ\u0081 -3.1017 -67.9356 2834.586
## 4 12240000                 PORTO SEGURO -3.3372 -67.4917 2789.684
## 5 12390000                 PORTO WALTER -8.2672 -72.7411 1095.260
## 6 12500000              CRUZEIRO DO SUL -7.6333 -72.6603 2078.195
\end{verbatim}

\begin{Shaded}
\begin{Highlighting}[]
\DocumentationTok{\#\# ============================================================}
\CommentTok{\# PARTE 6 — LER E JUNTAR TABELA DE RENDIMENTO HIDROLÓGICO}
\CommentTok{\# ============================================================}

\NormalTok{rendimento\_hidrologico }\OtherTok{\textless{}{-}} \FunctionTok{read\_csv2}\NormalTok{(}
\NormalTok{  tabela\_volume,}
  \AttributeTok{show\_col\_types =} \ConstantTok{FALSE}
\NormalTok{) }\SpecialCharTok{\%\textgreater{}\%}
  \FunctionTok{mutate}\NormalTok{(}
    \AttributeTok{Codigo =} \FunctionTok{as.character}\NormalTok{(Codigo)}
\NormalTok{  )}
\end{Highlighting}
\end{Shaded}

\begin{verbatim}
## i Using "','" as decimal and "'.'" as grouping mark. Use `read_delim()` for more control.
\end{verbatim}

\begin{Shaded}
\begin{Highlighting}[]
\FunctionTok{names}\NormalTok{(rendimento\_hidrologico)}
\end{Highlighting}
\end{Shaded}

\begin{verbatim}
##  [1] "Codigo"    "Nome"      "Latitude"  "Longitude" "ArDrngm"   "mm_2014"  
##  [7] "mm_2015"   "mm_2016"   "mm_2017"   "mm_2018"   "mm_2019"   "mm_2020"  
## [13] "mm_2021"   "mm_2022"   "mm_2023"   "mm_2024"   "mm_2025"   "mm_2026"
\end{verbatim}

\begin{Shaded}
\begin{Highlighting}[]
\NormalTok{tabela\_final }\OtherTok{\textless{}{-}}\NormalTok{ tabela\_final }\SpecialCharTok{\%\textgreater{}\%}
  \FunctionTok{left\_join}\NormalTok{(}
\NormalTok{    rendimento\_hidrologico,}
    \AttributeTok{by =} \StringTok{"Codigo"}
\NormalTok{  )}

\FunctionTok{head}\NormalTok{(tabela\_final)}
\end{Highlighting}
\end{Shaded}

\begin{verbatim}
##     Codigo                       Nome.x Latitud  Longitd    QF_v2
## 1 11200000                     TERESINA -4.3578 -69.7336 1461.153
## 2 11400000      SÃO PAULO DE OLIVENÇA -3.4569 -68.9119 2131.965
## 3 11500000 SANTO ANTÔNIO DO IÇÃ\u0081 -3.1017 -67.9356 2834.586
## 4 12240000                 PORTO SEGURO -3.3372 -67.4917 2789.684
## 5 12390000                 PORTO WALTER -8.2672 -72.7411 1095.260
## 6 12500000              CRUZEIRO DO SUL -7.6333 -72.6603 2078.195
##                  Nome.y Latitude Longitude ArDrngm  mm_2014   mm_2015   mm_2016
## 1              TERESINA  -4.3578  -69.7336  985000       NA        NA        NA
## 2 SÃO PAULO DE OLIVENÇA  -3.4569  -68.9119 1010000 1252.586 1757.1739 1380.7446
## 3  SANTO ANTÔNIO DO IÇÁ  -3.1017  -67.9356 1130000 1420.867 1995.4295 1569.8717
## 4          PORTO SEGURO  -3.3372  -67.4917   64400 1105.876 1494.5727 1059.5049
## 5          PORTO WALTER  -8.2672  -72.7411   22200    0.000  936.1247  768.6346
## 6       CRUZEIRO DO SUL  -7.6333  -72.6603   37800    0.000  818.4283  630.0941
##     mm_2017   mm_2018   mm_2019   mm_2020   mm_2021   mm_2022   mm_2023
## 1 1307.2894 1353.9380 1575.1049 1431.8784 1428.3081 1484.5948 1255.6737
## 2 1449.1055 1402.0214 1634.9759 1485.1973 1480.0249 1417.9907 1231.7547
## 3 1637.9528 1578.9300 1820.3814 1670.4325 1683.2525 1610.1537 1390.3726
## 4 1216.8473 1176.3035 1272.6270 1227.4590 1267.7182 1268.3292 1159.0364
## 5  768.4573  537.4223  850.0891  646.9406  731.6684  699.5538  539.0896
## 6  883.1662  530.0486 1006.8934  754.9957  847.2165  660.3205  626.9251
##     mm_2024   mm_2025   mm_2026
## 1 1108.0657 1486.9272 392.77740
## 2 1213.4510 1202.2990        NA
## 3 1313.7637 1359.2972        NA
## 4 1019.1373 1100.1150 125.37270
## 5  427.9440  488.2474  91.86335
## 6  716.3078  587.8716        NA
\end{verbatim}

\begin{Shaded}
\begin{Highlighting}[]
\FunctionTok{names}\NormalTok{(tabela\_final)}
\end{Highlighting}
\end{Shaded}

\begin{verbatim}
##  [1] "Codigo"    "Nome.x"    "Latitud"   "Longitd"   "QF_v2"     "Nome.y"   
##  [7] "Latitude"  "Longitude" "ArDrngm"   "mm_2014"   "mm_2015"   "mm_2016"  
## [13] "mm_2017"   "mm_2018"   "mm_2019"   "mm_2020"   "mm_2021"   "mm_2022"  
## [19] "mm_2023"   "mm_2024"   "mm_2025"   "mm_2026"
\end{verbatim}

\begin{Shaded}
\begin{Highlighting}[]
\CommentTok{\# ============================================================}
\CommentTok{\# PARTE 7 — CALCULAR MÉDIA DOS ANOS DE RENDIMENTO}
\CommentTok{\#            Selecionando os últimos 5 anos até o ano final}
\CommentTok{\# ============================================================}

\CommentTok{\# Escolha aqui o ano final e a quantidade de anos}
\NormalTok{ano\_final }\OtherTok{\textless{}{-}} \DecValTok{2024}
\NormalTok{n\_anos\_media }\OtherTok{\textless{}{-}} \DecValTok{5}

\NormalTok{anos\_media }\OtherTok{\textless{}{-}}\NormalTok{ (ano\_final }\SpecialCharTok{{-}}\NormalTok{ n\_anos\_media }\SpecialCharTok{+} \DecValTok{1}\NormalTok{)}\SpecialCharTok{:}\NormalTok{ano\_final}

\FunctionTok{cat}\NormalTok{(}\StringTok{"Anos usados na média:}\SpecialCharTok{\textbackslash{}n}\StringTok{"}\NormalTok{)}
\end{Highlighting}
\end{Shaded}

\begin{verbatim}
## Anos usados na média:
\end{verbatim}

\begin{Shaded}
\begin{Highlighting}[]
\FunctionTok{print}\NormalTok{(anos\_media)}
\end{Highlighting}
\end{Shaded}

\begin{verbatim}
## [1] 2020 2021 2022 2023 2024
\end{verbatim}

\begin{Shaded}
\begin{Highlighting}[]
\CommentTok{\# Criar nomes das colunas esperadas}
\NormalTok{colunas\_mm }\OtherTok{\textless{}{-}} \FunctionTok{paste0}\NormalTok{(}\StringTok{"mm\_"}\NormalTok{, anos\_media)}

\CommentTok{\# Conferir se essas colunas existem na tabela}
\NormalTok{colunas\_mm\_existentes }\OtherTok{\textless{}{-}}\NormalTok{ colunas\_mm[}
\NormalTok{  colunas\_mm }\SpecialCharTok{\%in\%} \FunctionTok{names}\NormalTok{(tabela\_final)}
\NormalTok{]}

\NormalTok{colunas\_mm\_faltando }\OtherTok{\textless{}{-}}\NormalTok{ colunas\_mm[}
  \SpecialCharTok{!}\NormalTok{colunas\_mm }\SpecialCharTok{\%in\%} \FunctionTok{names}\NormalTok{(tabela\_final)}
\NormalTok{]}

\FunctionTok{cat}\NormalTok{(}\StringTok{"Colunas encontradas:}\SpecialCharTok{\textbackslash{}n}\StringTok{"}\NormalTok{)}
\end{Highlighting}
\end{Shaded}

\begin{verbatim}
## Colunas encontradas:
\end{verbatim}

\begin{Shaded}
\begin{Highlighting}[]
\FunctionTok{print}\NormalTok{(colunas\_mm\_existentes)}
\end{Highlighting}
\end{Shaded}

\begin{verbatim}
## [1] "mm_2020" "mm_2021" "mm_2022" "mm_2023" "mm_2024"
\end{verbatim}

\begin{Shaded}
\begin{Highlighting}[]
\FunctionTok{cat}\NormalTok{(}\StringTok{"Colunas faltando:}\SpecialCharTok{\textbackslash{}n}\StringTok{"}\NormalTok{)}
\end{Highlighting}
\end{Shaded}

\begin{verbatim}
## Colunas faltando:
\end{verbatim}

\begin{Shaded}
\begin{Highlighting}[]
\FunctionTok{print}\NormalTok{(colunas\_mm\_faltando)}
\end{Highlighting}
\end{Shaded}

\begin{verbatim}
## character(0)
\end{verbatim}

\begin{Shaded}
\begin{Highlighting}[]
\ControlFlowTok{if}\NormalTok{ (}\FunctionTok{length}\NormalTok{(colunas\_mm\_existentes) }\SpecialCharTok{==} \DecValTok{0}\NormalTok{) \{}
  \FunctionTok{stop}\NormalTok{(}\StringTok{"Nenhuma coluna dos anos selecionados foi encontrada na tabela\_final."}\NormalTok{)}
\NormalTok{\}}

\CommentTok{\# Calcular média, mínimo, máximo e desvio padrão}
\CommentTok{\# NA é ignorado automaticamente}
\NormalTok{tabela\_final }\OtherTok{\textless{}{-}}\NormalTok{ tabela\_final }\SpecialCharTok{\%\textgreater{}\%}
  \FunctionTok{rowwise}\NormalTok{() }\SpecialCharTok{\%\textgreater{}\%}
  \FunctionTok{mutate}\NormalTok{(}
    \AttributeTok{n\_anos\_validos =} \FunctionTok{sum}\NormalTok{(}\SpecialCharTok{!}\FunctionTok{is.na}\NormalTok{(}\FunctionTok{c\_across}\NormalTok{(}\FunctionTok{all\_of}\NormalTok{(colunas\_mm\_existentes)))),}
    
    \AttributeTok{media\_mm =} \FunctionTok{ifelse}\NormalTok{(}
\NormalTok{      n\_anos\_validos }\SpecialCharTok{\textgreater{}} \DecValTok{0}\NormalTok{,}
      \FunctionTok{mean}\NormalTok{(}\FunctionTok{c\_across}\NormalTok{(}\FunctionTok{all\_of}\NormalTok{(colunas\_mm\_existentes)), }\AttributeTok{na.rm =} \ConstantTok{TRUE}\NormalTok{),}
      \ConstantTok{NA\_real\_}
\NormalTok{    ),}
    
    \AttributeTok{min\_mm =} \FunctionTok{ifelse}\NormalTok{(}
\NormalTok{      n\_anos\_validos }\SpecialCharTok{\textgreater{}} \DecValTok{0}\NormalTok{,}
      \FunctionTok{min}\NormalTok{(}\FunctionTok{c\_across}\NormalTok{(}\FunctionTok{all\_of}\NormalTok{(colunas\_mm\_existentes)), }\AttributeTok{na.rm =} \ConstantTok{TRUE}\NormalTok{),}
      \ConstantTok{NA\_real\_}
\NormalTok{    ),}
    
    \AttributeTok{max\_mm =} \FunctionTok{ifelse}\NormalTok{(}
\NormalTok{      n\_anos\_validos }\SpecialCharTok{\textgreater{}} \DecValTok{0}\NormalTok{,}
      \FunctionTok{max}\NormalTok{(}\FunctionTok{c\_across}\NormalTok{(}\FunctionTok{all\_of}\NormalTok{(colunas\_mm\_existentes)), }\AttributeTok{na.rm =} \ConstantTok{TRUE}\NormalTok{),}
      \ConstantTok{NA\_real\_}
\NormalTok{    ),}
    
    \AttributeTok{sd\_mm =} \FunctionTok{ifelse}\NormalTok{(}
\NormalTok{      n\_anos\_validos }\SpecialCharTok{\textgreater{}} \DecValTok{1}\NormalTok{,}
      \FunctionTok{sd}\NormalTok{(}\FunctionTok{c\_across}\NormalTok{(}\FunctionTok{all\_of}\NormalTok{(colunas\_mm\_existentes)), }\AttributeTok{na.rm =} \ConstantTok{TRUE}\NormalTok{),}
      \ConstantTok{NA\_real\_}
\NormalTok{    )}
\NormalTok{  ) }\SpecialCharTok{\%\textgreater{}\%}
  \FunctionTok{ungroup}\NormalTok{()}

\FunctionTok{head}\NormalTok{(}
\NormalTok{  tabela\_final }\SpecialCharTok{\%\textgreater{}\%}
    \FunctionTok{select}\NormalTok{(}
\NormalTok{      Codigo,}
      \FunctionTok{all\_of}\NormalTok{(colunas\_mm\_existentes),}
\NormalTok{      n\_anos\_validos,}
\NormalTok{      media\_mm,}
\NormalTok{      min\_mm,}
\NormalTok{      max\_mm,}
\NormalTok{      sd\_mm,}
\NormalTok{      QF\_v2}
\NormalTok{    )}
\NormalTok{)}
\end{Highlighting}
\end{Shaded}

\begin{verbatim}
## # A tibble: 6 x 12
##   Codigo  mm_2020 mm_2021 mm_2022 mm_2023 mm_2024 n_anos_validos media_mm min_mm
##   <chr>     <dbl>   <dbl>   <dbl>   <dbl>   <dbl>          <int>    <dbl>  <dbl>
## 1 112000~   1432.   1428.   1485.   1256.   1108.              5    1342.  1108.
## 2 114000~   1485.   1480.   1418.   1232.   1213.              5    1366.  1213.
## 3 115000~   1670.   1683.   1610.   1390.   1314.              5    1534.  1314.
## 4 122400~   1227.   1268.   1268.   1159.   1019.              5    1188.  1019.
## 5 123900~    647.    732.    700.    539.    428.              5     609.   428.
## 6 125000~    755.    847.    660.    627.    716.              5     721.   627.
## # i 3 more variables: max_mm <dbl>, sd_mm <dbl>, QF_v2 <dbl>
\end{verbatim}

\begin{Shaded}
\begin{Highlighting}[]
\CommentTok{\# ============================================================}
\CommentTok{\# PARTE 8 — GRÁFICO 1:1 QUADRADO E COM MESMA ESCALA}
\CommentTok{\# ============================================================}

\NormalTok{dados\_grafico }\OtherTok{\textless{}{-}}\NormalTok{ tabela\_final }\SpecialCharTok{\%\textgreater{}\%}
  \FunctionTok{filter}\NormalTok{(}
    \SpecialCharTok{!}\FunctionTok{is.na}\NormalTok{(media\_mm),}
    \SpecialCharTok{!}\FunctionTok{is.na}\NormalTok{(QF\_v2),}
    \FunctionTok{is.finite}\NormalTok{(media\_mm),}
    \FunctionTok{is.finite}\NormalTok{(QF\_v2)}
\NormalTok{  )}

\NormalTok{limites }\OtherTok{\textless{}{-}} \FunctionTok{range}\NormalTok{(}
  \FunctionTok{c}\NormalTok{(dados\_grafico}\SpecialCharTok{$}\NormalTok{media\_mm, dados\_grafico}\SpecialCharTok{$}\NormalTok{QF\_v2),}
  \AttributeTok{na.rm =} \ConstantTok{TRUE}
\NormalTok{)}

\NormalTok{margem }\OtherTok{\textless{}{-}} \FunctionTok{diff}\NormalTok{(limites) }\SpecialCharTok{*} \FloatTok{0.05}

\NormalTok{limites }\OtherTok{\textless{}{-}} \FunctionTok{c}\NormalTok{(}
\NormalTok{  limites[}\DecValTok{1}\NormalTok{] }\SpecialCharTok{{-}}\NormalTok{ margem,}
\NormalTok{  limites[}\DecValTok{2}\NormalTok{] }\SpecialCharTok{+}\NormalTok{ margem}
\NormalTok{)}

\NormalTok{grafico\_1\_1 }\OtherTok{\textless{}{-}} \FunctionTok{ggplot}\NormalTok{(}
\NormalTok{  dados\_grafico,}
  \FunctionTok{aes}\NormalTok{(}
    \AttributeTok{x =}\NormalTok{ media\_mm,}
    \AttributeTok{y =}\NormalTok{ QF\_v2}
\NormalTok{  )}
\NormalTok{) }\SpecialCharTok{+}
  \FunctionTok{geom\_point}\NormalTok{(}\AttributeTok{size =} \DecValTok{2}\NormalTok{, }\AttributeTok{alpha =} \FloatTok{0.75}\NormalTok{) }\SpecialCharTok{+}
  \FunctionTok{geom\_abline}\NormalTok{(}
    \AttributeTok{intercept =} \DecValTok{0}\NormalTok{,}
    \AttributeTok{slope =} \DecValTok{1}\NormalTok{,}
    \AttributeTok{linetype =} \StringTok{"dashed"}\NormalTok{,}
    \AttributeTok{linewidth =} \FloatTok{0.8}
\NormalTok{  ) }\SpecialCharTok{+}
  \FunctionTok{scale\_x\_continuous}\NormalTok{(}
    \AttributeTok{limits =}\NormalTok{ limites}
\NormalTok{  ) }\SpecialCharTok{+}
  \FunctionTok{scale\_y\_continuous}\NormalTok{(}
    \AttributeTok{limits =}\NormalTok{ limites}
\NormalTok{  ) }\SpecialCharTok{+}
  \FunctionTok{coord\_equal}\NormalTok{() }\SpecialCharTok{+}
  \FunctionTok{labs}\NormalTok{(}
    \AttributeTok{title =} \StringTok{"Comparação 1:1 do rendimento hídrico"}\NormalTok{,}
    \AttributeTok{subtitle =} \StringTok{"Média dos dados observados vs valor extraído do raster Quick flow"}\NormalTok{,}
    \AttributeTok{x =} \StringTok{"Rendimento hídrico observado médio (mm/ano)"}\NormalTok{,}
    \AttributeTok{y =} \StringTok{"Rendimento hídrico do raster Quick flow (mm/ano)"}
\NormalTok{  ) }\SpecialCharTok{+}
  \FunctionTok{theme\_bw}\NormalTok{()}

\NormalTok{grafico\_1\_1}
\end{Highlighting}
\end{Shaded}

\pandocbounded{\includegraphics[keepaspectratio]{4_comparative_vazao_invest_files/4_comparative_vazao_invest_files/figure-latex/unnamed-chunk-4-1.pdf}}

\begin{Shaded}
\begin{Highlighting}[]
\CommentTok{\# ============================================================}
\CommentTok{\# PARTE 9 — SALVAR CSV E GRÁFICO}
\CommentTok{\# ============================================================}

\FunctionTok{write\_csv2}\NormalTok{(}
\NormalTok{  tabela\_final,}
\NormalTok{  arquivo\_saida\_csv}
\NormalTok{)}

\NormalTok{arquivo\_grafico }\OtherTok{\textless{}{-}} \FunctionTok{file.path}\NormalTok{(}
\NormalTok{  dir\_saida,}
  \StringTok{"grafico\_1\_1\_rendimento\_hidrico.png"}
\NormalTok{)}

\FunctionTok{ggsave}\NormalTok{(}
  \AttributeTok{filename =}\NormalTok{ arquivo\_grafico,}
  \AttributeTok{plot =}\NormalTok{ grafico\_1\_1,}
  \AttributeTok{width =} \DecValTok{7}\NormalTok{,}
  \AttributeTok{height =} \DecValTok{7}\NormalTok{,}
  \AttributeTok{dpi =} \DecValTok{300}
\NormalTok{)}

\FunctionTok{cat}\NormalTok{(}\StringTok{"Tabela salva em:}\SpecialCharTok{\textbackslash{}n}\StringTok{"}\NormalTok{, arquivo\_saida\_csv, }\StringTok{"}\SpecialCharTok{\textbackslash{}n}\StringTok{"}\NormalTok{)}
\end{Highlighting}
\end{Shaded}

\begin{verbatim}
## Tabela salva em:
##  C:/trab/SWY_Amazonia/04_resultados/rendimento_hidrico_com_QF_v2.csv
\end{verbatim}

\begin{Shaded}
\begin{Highlighting}[]
\FunctionTok{cat}\NormalTok{(}\StringTok{"Gráfico salvo em:}\SpecialCharTok{\textbackslash{}n}\StringTok{"}\NormalTok{, arquivo\_grafico, }\StringTok{"}\SpecialCharTok{\textbackslash{}n}\StringTok{"}\NormalTok{)}
\end{Highlighting}
\end{Shaded}

\begin{verbatim}
## Gráfico salvo em:
##  C:/trab/SWY_Amazonia/04_resultados/grafico_1_1_rendimento_hidrico.png
\end{verbatim}

\begin{Shaded}
\begin{Highlighting}[]
\CommentTok{\# ============================================================}
\CommentTok{\# PARTE 10 — SALVAR SHAPEFILE FINAL COMPLETO}
\CommentTok{\# ============================================================}

\NormalTok{pontos\_shp\_final }\OtherTok{\textless{}{-}}\NormalTok{ pontos\_com\_id }\SpecialCharTok{\%\textgreater{}\%}
  \FunctionTok{left\_join}\NormalTok{(}
\NormalTok{    valores\_extraidos,}
    \AttributeTok{by =} \StringTok{"ID"}
\NormalTok{  ) }\SpecialCharTok{\%\textgreater{}\%}
  \FunctionTok{select}\NormalTok{(}
    \SpecialCharTok{{-}}\NormalTok{ID}
\NormalTok{  ) }\SpecialCharTok{\%\textgreater{}\%}
  \FunctionTok{left\_join}\NormalTok{(}
\NormalTok{    tabela\_final }\SpecialCharTok{\%\textgreater{}\%}
      \FunctionTok{st\_drop\_geometry}\NormalTok{() }\SpecialCharTok{\%\textgreater{}\%}
      \FunctionTok{select}\NormalTok{(}
\NormalTok{        Codigo,}
        \FunctionTok{all\_of}\NormalTok{(colunas\_mm),}
\NormalTok{        n\_anos\_validos,}
\NormalTok{        media\_mm,}
\NormalTok{        min\_mm,}
\NormalTok{        max\_mm,}
\NormalTok{        sd\_mm}
\NormalTok{      ),}
    \AttributeTok{by =} \StringTok{"Codigo"}
\NormalTok{  )}

\FunctionTok{st\_write}\NormalTok{(}
\NormalTok{  pontos\_shp\_final,}
\NormalTok{  arquivo\_saida\_shp,}
  \AttributeTok{delete\_layer =} \ConstantTok{TRUE}
\NormalTok{)}
\end{Highlighting}
\end{Shaded}

\begin{verbatim}
## Warning in abbreviate_shapefile_names(obj): Field names abbreviated for ESRI
## Shapefile driver
\end{verbatim}

\begin{verbatim}
## Deleting layer `rendimento_hidrico_com_QF_v2' using driver `ESRI Shapefile'
## Writing layer `rendimento_hidrico_com_QF_v2' to data source 
##   `C:/trab/SWY_Amazonia/04_resultados/rendimento_hidrico_com_QF_v2.shp' using driver `ESRI Shapefile'
## Writing 107 features with 15 fields and geometry type Point.
\end{verbatim}

\begin{verbatim}
## Warning in CPL_write_ogr(obj, dsn, layer, driver,
## as.character(dataset_options), : GDAL Message 1: One or several characters
## couldn't be converted correctly from UTF-8 to ISO-8859-1.  This warning will
## not be emitted anymore.
\end{verbatim}

\begin{Shaded}
\begin{Highlighting}[]
\FunctionTok{cat}\NormalTok{(}\StringTok{"Shapefile salvo em:}\SpecialCharTok{\textbackslash{}n}\StringTok{"}\NormalTok{, arquivo\_saida\_shp, }\StringTok{"}\SpecialCharTok{\textbackslash{}n}\StringTok{"}\NormalTok{)}
\end{Highlighting}
\end{Shaded}

\begin{verbatim}
## Shapefile salvo em:
##  C:/trab/SWY_Amazonia/04_resultados/rendimento_hidrico_com_QF_v2.shp
\end{verbatim}

\begin{Shaded}
\begin{Highlighting}[]
\CommentTok{\# ============================================================}
\CommentTok{\# TESTE — GRÁFICO 1:1 COM VARIAÇÃO ANUAL}
\CommentTok{\#          Ponto = média sem zeros}
\CommentTok{\#          Barra = mínimo{-}máximo sem zeros}
\CommentTok{\# ============================================================}

\NormalTok{dados\_grafico\_dispersao }\OtherTok{\textless{}{-}}\NormalTok{ tabela\_final }\SpecialCharTok{\%\textgreater{}\%}
  \FunctionTok{filter}\NormalTok{(}
    \SpecialCharTok{!}\FunctionTok{is.na}\NormalTok{(media\_mm),}
    \SpecialCharTok{!}\FunctionTok{is.na}\NormalTok{(QF\_v2),}
    \SpecialCharTok{!}\FunctionTok{is.na}\NormalTok{(min\_mm),}
    \SpecialCharTok{!}\FunctionTok{is.na}\NormalTok{(max\_mm),}
    \FunctionTok{is.finite}\NormalTok{(media\_mm),}
    \FunctionTok{is.finite}\NormalTok{(QF\_v2),}
    \FunctionTok{is.finite}\NormalTok{(min\_mm),}
    \FunctionTok{is.finite}\NormalTok{(max\_mm)}
\NormalTok{  )}

\NormalTok{limites\_teste }\OtherTok{\textless{}{-}} \FunctionTok{range}\NormalTok{(}
  \FunctionTok{c}\NormalTok{(}
\NormalTok{    dados\_grafico\_dispersao}\SpecialCharTok{$}\NormalTok{min\_mm,}
\NormalTok{    dados\_grafico\_dispersao}\SpecialCharTok{$}\NormalTok{max\_mm,}
\NormalTok{    dados\_grafico\_dispersao}\SpecialCharTok{$}\NormalTok{QF\_v2}
\NormalTok{  ),}
  \AttributeTok{na.rm =} \ConstantTok{TRUE}
\NormalTok{)}

\NormalTok{margem\_teste }\OtherTok{\textless{}{-}} \FunctionTok{diff}\NormalTok{(limites\_teste) }\SpecialCharTok{*} \FloatTok{0.05}

\NormalTok{limites\_teste }\OtherTok{\textless{}{-}} \FunctionTok{c}\NormalTok{(}
\NormalTok{  limites\_teste[}\DecValTok{1}\NormalTok{] }\SpecialCharTok{{-}}\NormalTok{ margem\_teste,}
\NormalTok{  limites\_teste[}\DecValTok{2}\NormalTok{] }\SpecialCharTok{+}\NormalTok{ margem\_teste}
\NormalTok{)}

\NormalTok{grafico\_1\_1\_dispersao\_teste }\OtherTok{\textless{}{-}} \FunctionTok{ggplot}\NormalTok{(}
\NormalTok{  dados\_grafico\_dispersao,}
  \FunctionTok{aes}\NormalTok{(}
    \AttributeTok{x =}\NormalTok{ media\_mm,}
    \AttributeTok{y =}\NormalTok{ QF\_v2}
\NormalTok{  )}
\NormalTok{) }\SpecialCharTok{+}
  \FunctionTok{geom\_segment}\NormalTok{(}
    \FunctionTok{aes}\NormalTok{(}
      \AttributeTok{x =}\NormalTok{ min\_mm,}
      \AttributeTok{xend =}\NormalTok{ max\_mm,}
      \AttributeTok{y =}\NormalTok{ QF\_v2,}
      \AttributeTok{yend =}\NormalTok{ QF\_v2}
\NormalTok{    ),}
    \AttributeTok{linewidth =} \FloatTok{0.6}\NormalTok{,}
    \AttributeTok{alpha =} \FloatTok{0.6}
\NormalTok{  ) }\SpecialCharTok{+}
  \FunctionTok{geom\_point}\NormalTok{(}
    \AttributeTok{size =} \FloatTok{2.5}\NormalTok{,}
    \AttributeTok{alpha =} \FloatTok{0.85}
\NormalTok{  ) }\SpecialCharTok{+}
  \FunctionTok{geom\_abline}\NormalTok{(}
    \AttributeTok{intercept =} \DecValTok{0}\NormalTok{,}
    \AttributeTok{slope =} \DecValTok{1}\NormalTok{,}
    \AttributeTok{linetype =} \StringTok{"dashed"}\NormalTok{,}
    \AttributeTok{linewidth =} \FloatTok{0.8}
\NormalTok{  ) }\SpecialCharTok{+}
  \FunctionTok{scale\_x\_continuous}\NormalTok{(}
    \AttributeTok{limits =}\NormalTok{ limites\_teste}
\NormalTok{  ) }\SpecialCharTok{+}
  \FunctionTok{scale\_y\_continuous}\NormalTok{(}
    \AttributeTok{limits =}\NormalTok{ limites\_teste}
\NormalTok{  ) }\SpecialCharTok{+}
  \FunctionTok{coord\_equal}\NormalTok{() }\SpecialCharTok{+}
  \FunctionTok{labs}\NormalTok{(}
    \AttributeTok{title =} \StringTok{"Comparação 1:1 com variação anual"}\NormalTok{,}
    \AttributeTok{subtitle =} \StringTok{"Ponto = média  | Barra = mínimo{-}máximo"}\NormalTok{,}
    \AttributeTok{x =} \StringTok{"Rendimento hídrico médio observado (mm/ano)"}\NormalTok{,}
    \AttributeTok{y =} \StringTok{"Rendimento hídrico do raster Quick Flow (mm/ano)"}
\NormalTok{  ) }\SpecialCharTok{+}
  \FunctionTok{theme\_bw}\NormalTok{()}

\NormalTok{grafico\_1\_1\_dispersao\_teste}
\end{Highlighting}
\end{Shaded}

\pandocbounded{\includegraphics[keepaspectratio]{4_comparative_vazao_invest_files/4_comparative_vazao_invest_files/figure-latex/unnamed-chunk-6-1.pdf}}

\begin{Shaded}
\begin{Highlighting}[]
\NormalTok{arquivo\_grafico\_teste }\OtherTok{\textless{}{-}} \FunctionTok{file.path}\NormalTok{(}
\NormalTok{  dir\_saida,}
  \StringTok{"teste\_grafico\_1\_1\_com\_variacao\_anual\_sem\_zeros.png"}
\NormalTok{)}

\FunctionTok{ggsave}\NormalTok{(}
  \AttributeTok{filename =}\NormalTok{ arquivo\_grafico\_teste,}
  \AttributeTok{plot =}\NormalTok{ grafico\_1\_1\_dispersao\_teste,}
  \AttributeTok{width =} \DecValTok{7}\NormalTok{,}
  \AttributeTok{height =} \DecValTok{7}\NormalTok{,}
  \AttributeTok{dpi =} \DecValTok{300}
\NormalTok{)}

\FunctionTok{cat}\NormalTok{(}\StringTok{"Gráfico salvo em:}\SpecialCharTok{\textbackslash{}n}\StringTok{"}\NormalTok{, arquivo\_grafico\_teste, }\StringTok{"}\SpecialCharTok{\textbackslash{}n}\StringTok{"}\NormalTok{)}
\end{Highlighting}
\end{Shaded}

\begin{verbatim}
## Gráfico salvo em:
##  C:/trab/SWY_Amazonia/04_resultados/teste_grafico_1_1_com_variacao_anual_sem_zeros.png
\end{verbatim}

\begin{Shaded}
\begin{Highlighting}[]
\CommentTok{\# ============================================================}
\CommentTok{\# LER TABELA DE VAZÕES}
\CommentTok{\# ============================================================}

\FunctionTok{library}\NormalTok{(dplyr)}
\FunctionTok{library}\NormalTok{(tidyr)}
\end{Highlighting}
\end{Shaded}

\begin{verbatim}
## 
## Anexando pacote: 'tidyr'
\end{verbatim}

\begin{verbatim}
## O seguinte objeto é mascarado por 'package:terra':
## 
##     extract
\end{verbatim}

\begin{Shaded}
\begin{Highlighting}[]
\FunctionTok{library}\NormalTok{(lubridate)}
\end{Highlighting}
\end{Shaded}

\begin{verbatim}
## 
## Anexando pacote: 'lubridate'
\end{verbatim}

\begin{verbatim}
## Os seguintes objetos são mascarados por 'package:terra':
## 
##     intersect, union
\end{verbatim}

\begin{verbatim}
## Os seguintes objetos são mascarados por 'package:base':
## 
##     date, intersect, setdiff, union
\end{verbatim}

\begin{Shaded}
\begin{Highlighting}[]
\FunctionTok{library}\NormalTok{(readr)}
\FunctionTok{library}\NormalTok{(ggplot2)}
\FunctionTok{library}\NormalTok{(stringr)}

\NormalTok{arquivo\_vazoes }\OtherTok{\textless{}{-}} \StringTok{"C:/trab/estacoes/tabela\_vazoes.csv"}

\NormalTok{tabela\_vazoes }\OtherTok{\textless{}{-}} \FunctionTok{read\_csv2}\NormalTok{(}
\NormalTok{  arquivo\_vazoes,}
  \AttributeTok{show\_col\_types =} \ConstantTok{FALSE}
\NormalTok{)}
\end{Highlighting}
\end{Shaded}

\begin{verbatim}
## i Using "','" as decimal and "'.'" as grouping mark. Use `read_delim()` for more control.
\end{verbatim}

\begin{Shaded}
\begin{Highlighting}[]
\NormalTok{tabela\_vazoes }\OtherTok{\textless{}{-}}\NormalTok{ tabela\_vazoes }\SpecialCharTok{\%\textgreater{}\%}
  \FunctionTok{mutate}\NormalTok{(}
    \AttributeTok{Data =} \FunctionTok{as.Date}\NormalTok{(Data)}
\NormalTok{  )}
\CommentTok{\# ============================================================}
\CommentTok{\# ORGANIZAR VAZÕES EM FORMATO LONGO}
\CommentTok{\# Excluir ano 2026}
\CommentTok{\# ============================================================}

\NormalTok{vazoes\_longas\_grafico }\OtherTok{\textless{}{-}}\NormalTok{ tabela\_vazoes }\SpecialCharTok{\%\textgreater{}\%}
  \FunctionTok{pivot\_longer}\NormalTok{(}
    \AttributeTok{cols =} \SpecialCharTok{{-}}\NormalTok{Data,}
    \AttributeTok{names\_to =} \StringTok{"Codigo"}\NormalTok{,}
    \AttributeTok{values\_to =} \StringTok{"Vazao"}
\NormalTok{  ) }\SpecialCharTok{\%\textgreater{}\%}
  \FunctionTok{mutate}\NormalTok{(}
    \AttributeTok{Ano =} \FunctionTok{year}\NormalTok{(Data),}
    \AttributeTok{Mes =} \FunctionTok{month}\NormalTok{(Data),}
    \AttributeTok{Vazao =} \FunctionTok{as.numeric}\NormalTok{(Vazao)}
\NormalTok{  ) }\SpecialCharTok{\%\textgreater{}\%}
  \FunctionTok{filter}\NormalTok{(}
\NormalTok{    Ano }\SpecialCharTok{!=} \DecValTok{2026}\NormalTok{,}
    \SpecialCharTok{!}\FunctionTok{is.na}\NormalTok{(Vazao)}
\NormalTok{  )}

\FunctionTok{head}\NormalTok{(vazoes\_longas\_grafico)}
\end{Highlighting}
\end{Shaded}

\begin{verbatim}
## # A tibble: 6 x 5
##   Data       Codigo    Vazao   Ano   Mes
##   <date>     <chr>     <dbl> <dbl> <dbl>
## 1 2014-01-01 11400000 45402.  2014     1
## 2 2014-01-01 11500000 54716.  2014     1
## 3 2014-01-01 12240000  2778.  2014     1
## 4 2014-01-01 12650000   430.  2014     1
## 5 2014-01-01 12840000  4128.  2014     1
## 6 2014-01-01 13150000 81651.  2014     1
\end{verbatim}

\begin{Shaded}
\begin{Highlighting}[]
\CommentTok{\# ============================================================}
\CommentTok{\# GRÁFICO 1 — VAZÃO MÉDIA ANUAL POR ESTAÇÃO}
\CommentTok{\# ============================================================}

\NormalTok{vazao\_anual }\OtherTok{\textless{}{-}}\NormalTok{ vazoes\_longas\_grafico }\SpecialCharTok{\%\textgreater{}\%}
  \FunctionTok{group\_by}\NormalTok{(}
\NormalTok{    Codigo,}
\NormalTok{    Ano}
\NormalTok{  ) }\SpecialCharTok{\%\textgreater{}\%}
  \FunctionTok{summarise}\NormalTok{(}
    \AttributeTok{vazao\_media\_anual =} \FunctionTok{mean}\NormalTok{(Vazao, }\AttributeTok{na.rm =} \ConstantTok{TRUE}\NormalTok{),}
    \AttributeTok{dias\_com\_dado =} \FunctionTok{n}\NormalTok{(),}
    \AttributeTok{.groups =} \StringTok{"drop"}
\NormalTok{  )}

\NormalTok{grafico\_vazao\_anual }\OtherTok{\textless{}{-}} \FunctionTok{ggplot}\NormalTok{(}
\NormalTok{  vazao\_anual,}
  \FunctionTok{aes}\NormalTok{(}
    \AttributeTok{x =}\NormalTok{ Ano,}
    \AttributeTok{y =}\NormalTok{ vazao\_media\_anual,}
    \AttributeTok{group =}\NormalTok{ Codigo}
\NormalTok{  )}
\NormalTok{) }\SpecialCharTok{+}
  \FunctionTok{geom\_line}\NormalTok{(}\AttributeTok{alpha =} \FloatTok{0.5}\NormalTok{) }\SpecialCharTok{+}
  \FunctionTok{geom\_point}\NormalTok{(}\AttributeTok{size =} \FloatTok{1.5}\NormalTok{, }\AttributeTok{alpha =} \FloatTok{0.7}\NormalTok{) }\SpecialCharTok{+}
  \FunctionTok{labs}\NormalTok{(}
    \AttributeTok{title =} \StringTok{"Comportamento anual da vazão por estação"}\NormalTok{,}
    \AttributeTok{subtitle =} \StringTok{"Vazão média anual — ano 2026 excluído"}\NormalTok{,}
    \AttributeTok{x =} \StringTok{"Ano"}\NormalTok{,}
    \AttributeTok{y =} \StringTok{"Vazão média anual (m³/s)"}
\NormalTok{  ) }\SpecialCharTok{+}
  \FunctionTok{theme\_bw}\NormalTok{()}

\NormalTok{grafico\_vazao\_anual}
\end{Highlighting}
\end{Shaded}

\pandocbounded{\includegraphics[keepaspectratio]{4_comparative_vazao_invest_files/4_comparative_vazao_invest_files/figure-latex/unnamed-chunk-7-1.pdf}}

\begin{Shaded}
\begin{Highlighting}[]
\FunctionTok{ggsave}\NormalTok{(}
  \AttributeTok{filename =} \StringTok{"C:/trab/estacoes/grafico\_vazao\_media\_anual.png"}\NormalTok{,}
  \AttributeTok{plot =}\NormalTok{ grafico\_vazao\_anual,}
  \AttributeTok{width =} \DecValTok{10}\NormalTok{,}
  \AttributeTok{height =} \DecValTok{6}\NormalTok{,}
  \AttributeTok{dpi =} \DecValTok{300}
\NormalTok{)}
\CommentTok{\# ============================================================}
\CommentTok{\# GRÁFICO 2 — MÉDIA GERAL ANUAL DAS ESTAÇÕES}
\CommentTok{\# ============================================================}

\NormalTok{vazao\_anual\_geral }\OtherTok{\textless{}{-}}\NormalTok{ vazao\_anual }\SpecialCharTok{\%\textgreater{}\%}
  \FunctionTok{group\_by}\NormalTok{(Ano) }\SpecialCharTok{\%\textgreater{}\%}
  \FunctionTok{summarise}\NormalTok{(}
    \AttributeTok{media\_geral =} \FunctionTok{mean}\NormalTok{(vazao\_media\_anual, }\AttributeTok{na.rm =} \ConstantTok{TRUE}\NormalTok{),}
    \AttributeTok{min\_geral =} \FunctionTok{min}\NormalTok{(vazao\_media\_anual, }\AttributeTok{na.rm =} \ConstantTok{TRUE}\NormalTok{),}
    \AttributeTok{max\_geral =} \FunctionTok{max}\NormalTok{(vazao\_media\_anual, }\AttributeTok{na.rm =} \ConstantTok{TRUE}\NormalTok{),}
    \AttributeTok{.groups =} \StringTok{"drop"}
\NormalTok{  )}

\NormalTok{grafico\_vazao\_anual\_geral }\OtherTok{\textless{}{-}} \FunctionTok{ggplot}\NormalTok{(}
\NormalTok{  vazao\_anual\_geral,}
  \FunctionTok{aes}\NormalTok{(}
    \AttributeTok{x =}\NormalTok{ Ano,}
    \AttributeTok{y =}\NormalTok{ media\_geral}
\NormalTok{  )}
\NormalTok{) }\SpecialCharTok{+}
  \FunctionTok{geom\_ribbon}\NormalTok{(}
    \FunctionTok{aes}\NormalTok{(}
      \AttributeTok{ymin =}\NormalTok{ min\_geral,}
      \AttributeTok{ymax =}\NormalTok{ max\_geral}
\NormalTok{    ),}
    \AttributeTok{alpha =} \FloatTok{0.2}
\NormalTok{  ) }\SpecialCharTok{+}
  \FunctionTok{geom\_line}\NormalTok{(}\AttributeTok{linewidth =} \DecValTok{1}\NormalTok{) }\SpecialCharTok{+}
  \FunctionTok{geom\_point}\NormalTok{(}\AttributeTok{size =} \DecValTok{2}\NormalTok{) }\SpecialCharTok{+}
  \FunctionTok{labs}\NormalTok{(}
    \AttributeTok{title =} \StringTok{"Comportamento anual médio da vazão"}\NormalTok{,}
    \AttributeTok{subtitle =} \StringTok{"Linha = média das estações | Faixa = mínimo{-}máximo entre estações | 2026 excluído"}\NormalTok{,}
    \AttributeTok{x =} \StringTok{"Ano"}\NormalTok{,}
    \AttributeTok{y =} \StringTok{"Vazão média anual (m³/s)"}
\NormalTok{  ) }\SpecialCharTok{+}
  \FunctionTok{theme\_bw}\NormalTok{()}

\NormalTok{grafico\_vazao\_anual\_geral}
\end{Highlighting}
\end{Shaded}

\pandocbounded{\includegraphics[keepaspectratio]{4_comparative_vazao_invest_files/4_comparative_vazao_invest_files/figure-latex/unnamed-chunk-7-2.pdf}}

\begin{Shaded}
\begin{Highlighting}[]
\FunctionTok{ggsave}\NormalTok{(}
  \AttributeTok{filename =} \StringTok{"C:/trab/estacoes/grafico\_vazao\_media\_anual\_geral.png"}\NormalTok{,}
  \AttributeTok{plot =}\NormalTok{ grafico\_vazao\_anual\_geral,}
  \AttributeTok{width =} \DecValTok{9}\NormalTok{,}
  \AttributeTok{height =} \DecValTok{6}\NormalTok{,}
  \AttributeTok{dpi =} \DecValTok{300}
\NormalTok{)}

\CommentTok{\# ============================================================}
\CommentTok{\# GRÁFICO 3 — DISTRIBUIÇÃO ANUAL DAS VAZÕES}
\CommentTok{\# ============================================================}

\NormalTok{grafico\_boxplot\_anual }\OtherTok{\textless{}{-}} \FunctionTok{ggplot}\NormalTok{(}
\NormalTok{  vazao\_anual,}
  \FunctionTok{aes}\NormalTok{(}
    \AttributeTok{x =} \FunctionTok{factor}\NormalTok{(Ano),}
    \AttributeTok{y =}\NormalTok{ vazao\_media\_anual}
\NormalTok{  )}
\NormalTok{) }\SpecialCharTok{+}
  \FunctionTok{geom\_boxplot}\NormalTok{() }\SpecialCharTok{+}
  \FunctionTok{labs}\NormalTok{(}
    \AttributeTok{title =} \StringTok{"Distribuição da vazão média anual entre estações"}\NormalTok{,}
    \AttributeTok{subtitle =} \StringTok{"Cada boxplot representa as estações em um ano | 2026 excluído"}\NormalTok{,}
    \AttributeTok{x =} \StringTok{"Ano"}\NormalTok{,}
    \AttributeTok{y =} \StringTok{"Vazão média anual (m³/s)"}
\NormalTok{  ) }\SpecialCharTok{+}
  \FunctionTok{theme\_bw}\NormalTok{()}

\NormalTok{grafico\_boxplot\_anual}
\end{Highlighting}
\end{Shaded}

\pandocbounded{\includegraphics[keepaspectratio]{4_comparative_vazao_invest_files/4_comparative_vazao_invest_files/figure-latex/unnamed-chunk-7-3.pdf}}

\begin{Shaded}
\begin{Highlighting}[]
\FunctionTok{ggsave}\NormalTok{(}
  \AttributeTok{filename =} \StringTok{"C:/trab/estacoes/grafico\_boxplot\_vazao\_anual.png"}\NormalTok{,}
  \AttributeTok{plot =}\NormalTok{ grafico\_boxplot\_anual,}
  \AttributeTok{width =} \DecValTok{10}\NormalTok{,}
  \AttributeTok{height =} \DecValTok{6}\NormalTok{,}
  \AttributeTok{dpi =} \DecValTok{300}
\NormalTok{)}
\end{Highlighting}
\end{Shaded}


\end{document}
